% 中文技术札记 —— 用 xelatex 编译（需要 ctex 与 Noto CJK 字体）
%   make            在本目录生成 main.pdf 并复制到 docs/
\documentclass[11pt,a4paper]{article}

\usepackage[UTF8,scheme=plain]{ctex}
\setCJKmainfont{Noto Serif CJK SC}
\setCJKsansfont{Noto Sans CJK SC}
\setCJKmonofont{Noto Sans Mono CJK SC}

\usepackage[margin=2.4cm]{geometry}
\usepackage{amsmath,amssymb,amsthm,mathtools}
\usepackage{booktabs}
\usepackage{graphicx}
\usepackage{xcolor}
\usepackage{caption}
\usepackage[colorlinks=true,allcolors=blue]{hyperref}
\usepackage{cleveref}
\usepackage{algorithm}
\usepackage{algpseudocode}
\usepackage{enumitem}

\crefname{section}{第}{第}
\crefformat{section}{第#2#1#3节}
\crefname{table}{表}{表}
\crefname{figure}{图}{图}
\crefname{equation}{式}{式}
\crefname{algorithm}{算法}{算法}

\newcommand{\Lop}{\mathcal{L}}
\newcommand{\Sop}{S_{\mathrm{op}}}
\newcommand{\bxi}{\boldsymbol{\xi}}
\newcommand{\NRMSE}{\mathrm{NRMSE}}

\title{\textbf{算子谱先验：受限传感器布局下的稀疏物理场重构}\\[4pt]
\large 技术札记（交接版）}
\author{}
\date{}

\begin{document}
\maketitle
\vspace{-2.4em}

\begin{abstract}
\noindent
在很多物理场重构问题里，传感器的位置由工程可达性决定，而不是由信息量决定：
只能贴在墙面、管壁或少数几个可达点上。此时重构不再是相邻观测之间的\textbf{内插}，
而是向从未观测过的区域\textbf{外推}，而平稳光滑核在这个任务上几乎无话可说——
它只知道"远处相关性低"，于是把远处预测成均值，那是关于我们无知的陈述，
不是关于场的陈述。

本札记给出一套从\textbf{方程形式}构造先验的方法：算子的符号给出解的联合功率谱，
把这个\textbf{不可分}的联合谱做\textbf{非负、保半正定}的低秩投影，
得到函数式张量分解每个 mode 的一维 GP 核。它只需要方程的形式和系数的一个范围，
\textbf{不需要任何调参数据}。

在三个算子族、二维与三维、九种传感器布局上：传感器散布时本方法与通用核\textbf{打平}
（这是应该的，内插不需要物理）；受限到单面墙时重构误差相对实际可部署的最好方法
下降 $22.8\%$ 至 $50.6\%$。我们同时如实报告方法失效的场景与两个尚未解释的现象。
\end{abstract}

\tableofcontents
\bigskip

%=====================================================================
%=====================================================================
\section{记号}
\label{sec:notation}

\begin{table}[h]
\centering
\small
\begin{tabular}{lll}
\toprule
符号 & 含义 & 形状 / 取值 \\
\midrule
$D$ & mode 数 & 二维场 $D=3$（$t,x,y$），三维场 $D=4$ \\
$n_d$ & 第 $d$ 个 mode 的格点数 & $64$（二维），$32$（三维） \\
$K_d$ & 第 $d$ 个 mode 的频率数 & $12$（二维），$8$（三维） \\
$R_d$ & 第 $d$ 个 mode 的 Tucker 秩 & $(8,5,5)$ / $(5,4,4,4)$ \\
$Q$ & bank 中的 atom 数 & $4$ \\
$\theta$ & 算子系数 $(D_x, D_y, r)$ & 真值 $\theta^\star$ 不可见 \\
\midrule
$z_d$ & 第 $d$ 个 mode 的连续坐标 & $z_d \in [0,1]$ \\
$\mathbf{i}(n)$ & 第 $n$ 个观测的多重索引 & $\mathbf{i}(n) \in \prod_d \{1,\dots,n_d\}$ \\
$\Omega$ & 观测索引集合 & $\lvert\Omega\rvert = 2621$（$1\%$） \\
$\bar{\Omega}$ & 留出索引集合 & 其余全部 \\
$\mathbf{y}$ & 带噪读数 & $\mathbb{R}^{\lvert\Omega\rvert}$ \\
\midrule
$\Sop$ & 算子给出的联合谱 & $K_1 \times \cdots \times K_D$ \\
$s_{dq}$ & 第 $d$ 个 mode 的第 $q$ 条谱 & $\mathbb{R}_{\ge 0}^{K_d}$ \\
$\pi_{dq}$ & 混合权重 & 单纯形，$\sum_q \pi_{dq}=1$ \\
$\bar{s}_d$ & 混合后的谱 $\sum_q \pi_{dq}s_{dq}$ & $\mathbb{R}_{\ge 0}^{K_d}$ \\
$\Phi_d$ & 本征基（无流余弦） & $n_d \times K_d$ \\
$\Psi_d$ & Mercer 特征 $\Phi_d \odot \sqrt{\bar{s}_d}$ & $n_d \times K_d$ \\
\midrule
$f^{(d)}_r$ & 第 $d$ 个 mode 第 $r$ 个因子函数 & GP \\
$\mathbf{a}^{(d)}_r$ & 该因子的展开系数 & $\mathbb{R}^{K_d}$ \\
$\mathbf{A}$ & 全部系数 $\{\mathbf{a}^{(d)}_r\}$ & $\sum_d R_d K_d = 216$ 个 \\
$\mathcal{G}$ & Tucker 核 & $R_1 \times \cdots \times R_D$ \\
$\sigma$ & 观测噪声标准差 & 标量，可学 \\
\midrule
$\mathbf{m}^{(d)}_r,\ \mathbf{v}^{(d)}_r$ & 变分后验的均值与方差 & $\mathbb{R}^{K_d}$ \\
\bottomrule
\end{tabular}
\end{table}


\section{引言与动机}

\subsection{一个具体场景}

一个厂房里有气体泄漏。你想知道\textbf{整个房间}的浓度分布——哪里最浓、
往哪个方向扩、要不要疏散。但你能装传感器的地方由不得你选：

\begin{itemize}[nosep,leftmargin=1.6em]
  \item 房间中央是设备和通道，不能布线也不能悬挂；
  \item 能开孔的只有墙面，而且往往只有可达的那一面；
  \item 燃烧室只能在壁面开测点，管道监测只能贴外壁，
        地下水污染监测只能打有限几口井。
\end{itemize}

于是观测集中在域的一小块\textbf{连续且位置很差}的区域里，
而需要重建的是\textbf{其余全部}。这个约束在物理场重构中是常态而非例外。

\subsection{为什么这不是标准的张量补全设定}

张量补全文献默认的缺失模式是\textbf{均匀随机}。两者的差别不是程度问题，
是性质问题（\cref{tab:setting}）。

\begin{table}[t]
\centering
\caption{两种缺失模式的差别。最后一行是本文的实测。}
\label{tab:setting}
\begin{tabular}{lll}
\toprule
 & 随机缺失 & 受限布局 \\
\midrule
未观测点有没有观测过的邻居 & 总是有 & 大部分一个都没有 \\
任务本质 & 内插 & \textbf{外推} \\
光滑先验够不够 & 够 & 不够 \\
实测（$1\%$ 观测） & 三种方法打平在 $0.053$ & 通用方法崩到 $0.67$--$0.96$ \\
\bottomrule
\end{tabular}
\end{table}

最后一行值得强调：\textbf{随机掩码下本方法与通用核完全打平}。
这不是失败，而是应该的——那个设定里物理确实没有可加的东西。

\subsection{光滑先验在外推时说了什么}

平稳核（Matérn / RBF）只说一句话：\textbf{相关性随距离衰减}。

内插时这句话够用：目标点旁边就有观测，"和邻居相似"直接给出答案。
外推时，同一句话变成：离观测区一远，相关性趋于零，\textbf{所以场趋于其均值}。

这是关于\emph{我们的无知}的陈述，不是关于\emph{场}的陈述。
它没说错，但它什么也没说。而且这件事能直接看见：
用传感器自己的读数调出来的 Matérn，重建出墙边条带之后，整个房间塌回零
（\cref{fig:recon} 第三列）。

\subsection{物理提供的恰好是缺的那句话}

关键观察是：\textbf{方程的形式几乎总是已知的}，哪怕系数不知道、解更不知道。
工程师往往不知道场长什么样，却知道自己面对的是扩散、输运还是波动。

对耗散算子，方程直接给出场沿每个轴如何衰减：

\begin{equation}
\Lop u = w \quad\Longrightarrow\quad
S_u(\bxi) = \bigl\lvert \hat{\Lop}(\bxi) \bigr\rvert^{-2} S_w(\bxi).
\label{eq:transfer}
\end{equation}

一句话概括二者的差别：

\begin{quote}
光滑核只知道"远处相关性低"；算子知道"\textbf{远处的值应该是多少}"。
\end{quote}

举一个能直接对上的例子：各向异性扩散的符号是
$i\omega + r + D_x k_x^2 + D_y k_y^2$。$D_x \neq D_y$ 意味着两个空间方向的
\textbf{衰减律不同}——沿一个轴平滑、沿另一个轴粗糙。
一个共享单一长度尺度的通用核没有办法表达这件事，而方程免费给出它。

\subsection{而且在这个设定里连"调参"都做不到}

最直接的反驳是："通用核也有超参，调一调不就行了？"
在受限布局下调不了，原因是结构性的。

调参需要\textbf{与预测目标相像的验证数据}。传感器全在一面墙上时，
你能留出来做验证的每一个点\emph{也在那面墙上}。
于是验证衡量的是\textbf{条带内部的内插}，而部署要求的是\textbf{横跨整个房间的外推}，
而这个划分看不到这个差别。实测见 \cref{tab:tuning}。

\begin{table}[t]
\centering
\caption{验证选出的长度尺度与实际需要的值。同一网格、同一宿主，只差谁来选。}
\label{tab:tuning}
\begin{tabular}{lll}
\toprule
布局 & 验证（传感器读数）选出 & 留出区域实际需要 \\
\midrule
房间内任意 & $0.8$ & $0.8$ \\
单面墙 & $0.5,\ 0.5,\ 0.8$ & $\mathbf{1.6,\ 2.4,\ 2.4}$ \\
三维房间单面墙 & $0.12$ & $\mathbf{3.5}$ \\
\bottomrule
\end{tabular}
\end{table}

散布传感器时验证选得准；受限时它一致地选出短 $3$--$5$ 倍
（三维上短 $30$ 倍）的尺度。这不是噪声，是系统性的——它衡量的本来就是另一件事。
\textbf{而方程直接给出核，不需要选。}

%=====================================================================
\section{背景}

\subsection{函数式张量分解}

把物理场离散成张量 $\mathcal{X} \in \mathbb{R}^{n_1 \times \cdots \times n_D}$，
低秩分解利用多向相关性。\emph{函数式}版本把离散因子表换成连续函数，
于是未观测坐标可以被插值而不是查表。以 Tucker 为例：

\begin{equation}
u(z_1, \dots, z_D) \;\approx\; \sum_{r_1, \dots, r_D}
  \mathcal{G}_{r_1 \cdots r_D} \prod_{d=1}^{D} f^{(d)}_{r_d}(z_d),
\label{eq:tucker}
\end{equation}

其中每个 $f^{(d)}_{r}$ 是一个一维 GP。核心问题随之而来：
\textbf{这些 GP 的核从哪来？}标准做法是给每个 mode 配一个通用 Matérn / RBF 核，
或者从同一批稀缺数据里学核。本文改变的正是这一点，其余全部保持不变。

\subsection{为什么宿主必须是 Tucker 而不是 CP}

真实二维场是一堆 blob，不是外积。它的 multilinear rank 小，但 CP rank 不小——
在稀疏观测能识别的秩下根本表示不了。CP 宿主下任何核都帮不上忙，
因为模型压根表示不了这个场。

Tucker 还允许\textbf{逐 mode 独立的秩}，这不是美观而是必要：
multilinear rank 为 $[2, 11, 11]$ 的场需要 $2 \cdot 11 \cdot 11 = 242$ 的核，
而强行取等秩 $11$ 需要 $1331$——这个差别决定模型在真实观测量下是否可辨识。

%=====================================================================
\section{问题设定}
\label{sec:setting}

给定：

\begin{itemize}[nosep,leftmargin=1.6em]
  \item 一个真实场 $\mathcal{X}$，由算子 $\Lop_{\theta^\star}$ 驱动，$\theta^\star$ 未知；
  \item 一个观测索引集合 $\Omega$，其位置由\textbf{可达性}决定（贴墙、贴面、局部块），
        而不是随机采样；
  \item 带噪读数 $y_i = \mathcal{X}[\omega_i] + \varepsilon_i$，$\omega_i \in \Omega$；
  \item 算子的\textbf{族}（扩散 / 输运 / 波动）已知，
        系数只知道落在一个预先声明的范围 $\Theta$ 内。
\end{itemize}

目标：在留出集合 $\bar{\Omega} = \text{全体} \setminus \Omega$ 上最小化

\begin{equation}
\NRMSE = \frac{\bigl\lVert \hat{u}[\bar{\Omega}] - \mathcal{X}[\bar{\Omega}] \bigr\rVert_2}
              {\sqrt{\lvert \bar{\Omega} \rvert}\ \operatorname{std}\bigl(\mathcal{X}[\bar{\Omega}]\bigr)} .
\label{eq:nrmse}
\end{equation}

按留出集标准差归一化，所以 $\NRMSE = 1.0$ \textbf{恰好是"预测均值"的分数}。
这让每个数字不需要参照任何 baseline 就能读懂。

\textbf{关键约束}：不允许使用任何用于选超参的额外数据。
这不是自我设限，而是\cref{tab:tuning}所示的现实——受限布局下那种数据不存在。

%=====================================================================
\section{方法}

方法分三步：从算子得到联合谱；把不可分的联合谱投影成每维一条谱；
把每维谱变成合法的核并放进宿主。

\subsection{第一步：从算子到解的功率谱}

设场 $u$ 由常系数线性算子 $\Lop$ 驱动，受随机强迫 $w$，即 $\Lop u = w$。
做时空 Fourier 变换，记 $\hat{\Lop}(\bxi)$ 为算子的\textbf{符号}，
$\bxi = (\omega, k_x, k_y)$。线性系统的输出谱由\cref{eq:transfer}给出。
取 $w$ 为白噪（$S_w \equiv \sigma^2$）：

\begin{equation}
\Sop(\bxi) \;\propto\; \bigl\lvert \hat{\Lop}(\bxi) \bigr\rvert^{-2}.
\end{equation}

主线算子是各向异性反应-扩散
$\partial_t u = D_x \partial_{xx} u + D_y \partial_{yy} u - r u + w$，
符号为 $\hat{\Lop} = i\omega + r + D_x k_x^2 + D_y k_y^2$，于是

\begin{equation}
\Sop(\omega, k_x, k_y) \;=\; \frac{1}{\omega^2 + \bigl(r + D_x k_x^2 + D_y k_y^2\bigr)^2}.
\label{eq:joint}
\end{equation}

\paragraph{必须用离散算子的特征值。}
在 $n$ 个格点、无流边界下，二阶差分算子的特征值不是 $k^2$ 而是

\begin{equation}
\lambda(k) \;=\; \frac{2 - 2\cos(\pi k / n)}{h^2}, \qquad h = 1/n .
\label{eq:eig}
\end{equation}

这一步曾经出过大错：早期版本用手写的 index-space 系数 $(1.0, 0.3)$，
而物理系数是 $(0.02, 0.006)$，导致先验衰减快了约 $4$ 倍。
\textbf{必须从离散算子的特征值构造，不能从连续 $k^2$ 直接套。}

\subsection{第二步：不可分性——本方法要解决的核心问题}

$\Sop$ 是一个 $D$ 阶张量，而\cref{eq:tucker}的宿主要的是\textbf{每个 mode 一条一维谱}。
二者不兼容，因为对任何一组一维函数都有

\begin{equation}
\Sop(\omega, k_x, k_y) \;\neq\; S_t(\omega)\, S_x(k_x)\, S_y(k_y),
\end{equation}

\cref{eq:joint}分母里的 $\bigl(r + D_x k_x^2 + D_y k_y^2\bigr)^2$ \textbf{天然不可分}。

于是问题变成：\emph{如何把一个不可分的联合谱投影成可分的、
且仍然对应合法核的每维谱？}

\subsection{第三步：非负 CP 投影}

用秩 $Q$ 的\textbf{非负} CP 分解逼近联合谱：

\begin{equation}
\Sop \;\approx\; \sum_{q=1}^{Q} c_q \; s^{(t)}_q \otimes s^{(x)}_q \otimes s^{(y)}_q,
\qquad s^{(d)}_q \ge 0 .
\label{eq:cp}
\end{equation}

\textbf{非负是必须的而非美观}：谱必须非负才对应一个半正定核。
若用普通 CP，因子可以取负，得到的"核"不是核。

求解用乘性更新（Euclidean loss），第 $d$ 个 mode 的更新为

\begin{equation}
s^{(d)} \;\leftarrow\; s^{(d)} \odot
  \frac{X_{(d)} \, K_{(d)}}{s^{(d)} \bigl(K_{(d)}^\top K_{(d)}\bigr) + \varepsilon},
\end{equation}

其中 $X_{(d)}$ 是沿 mode $d$ 的展开矩阵，$K_{(d)}$ 是其余因子的 Khatri--Rao 积。

\paragraph{掩码而不是置零。}
数据做了去均值，这恰好抹掉了\textbf{联合模态} $(0,0,0)$。
早期实现是把该元素\emph{置零}再分离，这是个真实的缺陷：
一个秩 $Q$ 的\emph{可分}模型无法只在一个角上取零，
它只能通过压掉\emph{某个因子}的 $k=0$ 来实现，而它会压掉最"便宜"的那个。
实测压掉的是时间维——真实场沿时间近乎常数（$k=0$ 占 $0.837$ 的能量），
而置零版先验只给了 $0.107$，把 $0.514$ 压在 $k=1$ 上，\textbf{先验在和数据对着干}。

正确做法是让 CP \textbf{忽略}该元素。带掩码 $M$ 时 Gram 捷径失效，必须显式加权：

\begin{equation}
s^{(d)} \;\leftarrow\; s^{(d)} \odot
  \frac{\bigl(M_{(d)} \odot X_{(d)}\bigr) K_{(d)}}
       {\bigl(M_{(d)} \odot (s^{(d)} K_{(d)}^\top)\bigr) K_{(d)} + \varepsilon}.
\end{equation}

逐模态 KL（与真实场的经验谱相比）：时间维 $1.578 \to 0.455$，
总和 $2.358 \to 1.711$（对照：调好的 Matérn 是 $4.204$）。

\subsection{第四步：从谱到核（Mercer）}

给定一维谱 $s_d$ 和一组正交基 $\{\phi_{dq}\}$，核由 Mercer 展开给出：

\begin{equation}
k_d(z, z') \;=\; \sum_{q} s_d(q) \, \phi_{dq}(z) \, \phi_{dq}(z') .
\end{equation}

关键点：当 $\{\phi_{dq}\}$ \textbf{就是算子的本征函数}时，
这个有限展开不是近似而是\textbf{精确的} Mercer 特征映射。
所以这里用 inducing points 反而是\textbf{严格的降级}——
它会用低秩近似去逼近一个本来就精确的有限秩表示。

实现上取特征

\begin{equation}
\Phi_d(z) \;=\; \bigl[\, \sqrt{s_d(0)}\,\phi_{d0}(z), \; \ldots, \;
                        \sqrt{s_d(Q)}\,\phi_{dQ}(z) \,\bigr],
\end{equation}

于是 $k_d(z,z') = \Phi_d(z)^\top \Phi_d(z')$ \textbf{自动半正定}。

\subsection{第五步：边界条件决定本征基}

算子的本征基由\textbf{边界条件}决定，不是由符号决定：

\begin{itemize}[nosep,leftmargin=1.6em]
  \item 周期边界 $\Rightarrow$ 复指数 $e^{i 2\pi k x}$；
  \item \textbf{无流（Neumann）边界} $\Rightarrow$ 余弦
        $\phi_k(x) = \sqrt{2}\cos(\pi k x)$，$\phi_0 \equiv 1$。
\end{itemize}

房间的墙是无流的，初值数据沿时间也不是周期的；
强行施加 $f(0) = f(1)$ 是一个很大的非物理约束。
\textbf{实测}（单面墙，$3$ seeds，只换基）：周期 Fourier $0.5781$ 对无流余弦 $0.5448$。

归一化也随之改变：周期基下平均逐点方差是 $s_0 + 2\sum_{k>0} s_k$；
余弦基下每个特征的均方都是 $1$，所以是 $\sum_k s_k$。用错会引入系统性的整体尺度偏差。

\subsection{第六步：混合权重与知识档位}

每个 mode 的谱是 bank 内 atom 的凸组合，权重 $\pi$ 由 softmax 参数化并从数据学：

\begin{equation}
s_d \;=\; \sum_{q} \pi_{dq} \, s_{dq}, \qquad
\pi_{dq} \ge 0, \quad \textstyle\sum_q \pi_{dq} = 1 .
\end{equation}

\textbf{知识档位}由 bank 怎么造来定义（\cref{tab:ladder-def}）。
在 K1 / K0 下，$\pi$ 实际在做\textbf{软参数推断}——
它在"物理可达的谱"张成的集合内选择，而不是在无约束核空间里搜索。

\begin{table}[t]
\centering
\caption{知识档位。这是本构造与 PINN / AutoIP 差别\emph{在种类上}的地方：
后两者必须先承诺一个具体的 $\theta$ 才能写出残差或虚拟观测。}
\label{tab:ladder-def}
\begin{tabular}{lll}
\toprule
档 & 学习者知道什么 & bank 构造 \\
\midrule
K2 & 真值 $\theta^\star$ & 分离 $\Sop(\theta^\star)$，得 $Q$ 个 atom \\
K1 & $\theta^\star \in \Theta_1$（预先声明的范围） & 从 $\Theta_1$ 采 $M$ 组，各自分离，atom 池化 \\
K0 & $\theta^\star \in \Theta_0 \supset \Theta_1$ & 同上，范围更宽 \\
K$-$1 & 无算子信息 & 通用字典 \\
\bottomrule
\end{tabular}
\end{table}

\subsection{贝叶斯模型：联合概率}
\label{sec:joint}

前面几步给出的是\emph{先验}。这一节写清楚整个模型的联合概率、
推断怎么做、以及后验到底长什么样。

\paragraph{生成过程。}模型是

\begin{align}
\mathbf{a}^{(d)}_r &\sim \mathcal{N}(\mathbf{0}, \mathbf{I}_{K_d}),
  && d = 1,\dots,D,\quad r = 1,\dots,R_d, \label{eq:prior-a}\\
f^{(d)}_r(z_d) &= \Psi_d(z_d)^\top \mathbf{a}^{(d)}_r,
  && \Psi_d(z_d) = \Phi_d(z_d) \odot \sqrt{\bar{s}_d}, \label{eq:factor}\\
u(\mathbf{z}) &= \sum_{r_1=1}^{R_1}\cdots\sum_{r_D=1}^{R_D}
  \mathcal{G}_{r_1\cdots r_D} \prod_{d=1}^{D} f^{(d)}_{r_d}(z_d), \label{eq:field}\\
y_n \mid u &\sim \mathcal{N}\bigl(u(\mathbf{z}_{\mathbf{i}(n)}),\ \sigma^2\bigr),
  && n = 1,\dots,\lvert\Omega\rvert. \label{eq:lik}
\end{align}

\paragraph{为什么\cref{eq:prior-a}--\cref{eq:factor}就是那个 GP。}
把标准正态系数代入\cref{eq:factor}，因子函数的协方差是

\begin{equation}
\operatorname{Cov}\bigl[f^{(d)}_r(z),\, f^{(d)}_r(z')\bigr]
  = \Psi_d(z)^\top \mathbb{E}\bigl[\mathbf{a}\mathbf{a}^\top\bigr] \Psi_d(z')
  = \Psi_d(z)^\top \Psi_d(z')
  = \sum_{k} \bar{s}_d(k)\,\phi_{dk}(z)\,\phi_{dk}(z')
  = k_d(z, z'),
\end{equation}

正是第四步的 Mercer 展开。\textbf{关键在于这里没有近似}：
当 $\{\phi_{dk}\}$ 是算子的本征函数时，有限展开\emph{精确}给出该核，
所以\cref{eq:factor}是一个 GP 的精确重参数化，而不是一个低秩近似。
这也是为什么在这里用 inducing points 是严格的降级。

于是联合概率为

\begin{equation}
p(\mathbf{y}, \mathbf{A} \mid \mathcal{G}, \boldsymbol{\pi}, \sigma)
= \underbrace{\prod_{n=1}^{\lvert\Omega\rvert}
    \mathcal{N}\bigl(y_n \,\big|\, u(\mathbf{z}_{\mathbf{i}(n)}; \mathbf{A}, \mathcal{G}, \boldsymbol{\pi}),\ \sigma^2\bigr)}_{\text{似然}}
  \; \underbrace{\prod_{d=1}^{D}\prod_{r=1}^{R_d}
    \mathcal{N}\bigl(\mathbf{a}^{(d)}_r \,\big|\, \mathbf{0}, \mathbf{I}\bigr)}_{\text{算子给出的先验}} .
\label{eq:joint}
\end{equation}

\paragraph{哪些量是贝叶斯的，哪些不是。}这一点必须说清楚，
因为\cref{eq:joint}的写法容易让人以为一切都被积掉了：

\begin{itemize}[nosep,leftmargin=1.6em]
  \item \textbf{系数 $\mathbf{A}$ 是贝叶斯的}——有先验、有变分后验、被积分掉；
  \item \textbf{Tucker 核 $\mathcal{G}$、混合权重 $\boldsymbol{\pi}$、噪声 $\sigma$ 是点估计}，
        与变分参数一起用梯度优化（即 II 型极大似然 / 经验贝叶斯）。
\end{itemize}

这是刻意的取舍：$\mathcal{G}$ 有 $\prod_d R_d = 200$ 个自由度，
在 $1\%$ 观测下对它做完整推断既不稳定也不必要；
而混合权重 $\boldsymbol{\pi}$ 只有 $Q = 4$ 个，
把它当点估计等于在"物理可达的谱"张成的集合内做一次极大似然选择——
这正是\cref{tab:ladder-def}中 K1 / K0 档"软参数推断"的含义。

\paragraph{注意$\,$\cref{eq:field}对 $\mathbf{A}$ 不是线性的。}
似然对每个单独的 $\mathbf{a}^{(d)}_r$ 是线性高斯的，
但 $D$ 个 mode 的\emph{乘积}使整体成为多线性模型，
所以后验没有闭式，必须做近似推断。

\subsection{变分推断}

取对角高斯的平均场族：

\begin{equation}
q(\mathbf{A}) = \prod_{d=1}^{D}\prod_{r=1}^{R_d}
  \mathcal{N}\bigl(\mathbf{a}^{(d)}_r \,\big|\, \mathbf{m}^{(d)}_r,\ \operatorname{diag}(\mathbf{v}^{(d)}_r)\bigr).
\label{eq:q}
\end{equation}

证据下界为

\begin{equation}
\mathcal{F}(q, \mathcal{G}, \boldsymbol{\pi}, \sigma)
= \underbrace{\mathbb{E}_{q}\bigl[\log p(\mathbf{y} \mid \mathbf{A}, \mathcal{G}, \boldsymbol{\pi}, \sigma)\bigr]}_{\text{期望对数似然}}
  \;-\; \underbrace{\mathrm{KL}\bigl(q(\mathbf{A}) \,\Vert\, p(\mathbf{A})\bigr)}_{\text{有闭式}} .
\label{eq:elbo}
\end{equation}

\paragraph{KL 项有闭式。}因为先验\cref{eq:prior-a}是标准正态
（尺度已经吸进了特征 $\Psi_d$ 的 $\sqrt{\bar{s}_d}$），

\begin{equation}
\mathrm{KL}\bigl(q \Vert p\bigr) = \frac{1}{2}\sum_{d,r,k}
  \Bigl[\,\bigl(m^{(d)}_{rk}\bigr)^2 + v^{(d)}_{rk} - 1 - \log v^{(d)}_{rk}\,\Bigr].
\end{equation}

\textbf{这是把谱吸进特征的实际好处}：先验永远是 $\mathcal{N}(\mathbf{0},\mathbf{I})$，
不管算子给出的谱是什么样，KL 都不需要采样。

\paragraph{期望对数似然用重参数化采样。}由于\cref{eq:field}的乘积结构，
该期望没有闭式。取 $S = 3$ 个样本，
$\tilde{\mathbf{a}}^{(d),s}_r = \mathbf{m}^{(d)}_r + \sqrt{\mathbf{v}^{(d)}_r} \odot \boldsymbol{\epsilon}^{s}$，
$\boldsymbol{\epsilon}^{s} \sim \mathcal{N}(\mathbf{0},\mathbf{I})$：

\begin{equation}
\mathbb{E}_{q}\bigl[\log p(\mathbf{y} \mid \cdot)\bigr]
\;\approx\; \frac{1}{S}\sum_{s=1}^{S}\sum_{n \in \mathcal{B}}
  \log \mathcal{N}\bigl(y_n \,\big|\, \hat{u}^{s}(\mathbf{z}_{\mathbf{i}(n)}),\ \sigma^2\bigr)
  \cdot \frac{\lvert\Omega\rvert}{\lvert\mathcal{B}\rvert},
\end{equation}

其中 $\mathcal{B}$ 是 minibatch，末项是无偏放大因子。
\textbf{只在 $\Omega$ 上求值}，所以 $64^3$ 的张量在训练中从不物化。

\subsection{后验长什么样}

\paragraph{后验均值（本文报告的重构）。}
实现取\textbf{插值估计}：把变分均值直接代入\cref{eq:field}，

\begin{equation}
\hat{u}(\mathbf{z}) = \sum_{r_1 \cdots r_D} \mathcal{G}_{r_1\cdots r_D}
  \prod_{d=1}^{D} \Psi_d(z_d)^\top \mathbf{m}^{(d)}_{r_d}.
\label{eq:postmean}
\end{equation}

\textbf{这不是乘积的真后验均值。}对独立随机变量
$\mathbb{E}\bigl[\prod_d f^{(d)}\bigr] = \prod_d \mathbb{E}\bigl[f^{(d)}\bigr]$ 成立，
所以在平均场假设下\cref{eq:postmean}\emph{是}无偏的；
但平均场本身忽略了 mode 之间的后验相关性，
所以它是对真后验均值的近似，而不是它本身。本文全部 $\NRMSE$ 都基于\cref{eq:postmean}。

\paragraph{后验预测分布。}对新坐标 $\mathbf{z}^\ast$，

\begin{equation}
p(u^\ast \mid \mathbf{y}) \;\approx\; \int \delta\bigl(u^\ast - u(\mathbf{z}^\ast; \mathbf{A})\bigr)\, q(\mathbf{A})\, \mathrm{d}\mathbf{A},
\end{equation}

用蒙特卡洛近似：抽 $\tilde{\mathbf{A}}^{s} \sim q$，得到样本
$u^{\ast,s} = u(\mathbf{z}^\ast; \tilde{\mathbf{A}}^{s})$；
若要含观测噪声再加 $\mathcal{N}(0,\sigma^2)$。
在 CP 宿主下还有一条闭式的矩传播路径：对独立因子，

\begin{equation}
\operatorname{Var}\Bigl[\prod_d f^{(d)}_r\Bigr]
  = \prod_d \Bigl(\bigl(m_r^{(d)}\bigr)^2 + v_r^{(d)}\Bigr) - \prod_d \bigl(m_r^{(d)}\bigr)^2 ,
\end{equation}

Tucker 宿主因为核不是对角的，统一走采样。

\paragraph{不确定性是欠覆盖的，必须如实说。}
在 $1\%$ 观测下，名义 $95\%$ 的预测区间实际覆盖 $83.7\%$
（通用先验是 $77.5\%$），到 $5\%$ 观测才到 $91\%$。
本方法的区间也更宽（$1.14$ 对 $1.00$），所以它相对通用先验的优势
有一部分是"对不确定性更诚实"而不是"更准"。
\textbf{本文主张的是更好的点预测先验，不是更好的不确定性量化}——
这一条对任何想在此基础上做贝叶斯决策（例如源定位的置信区域）的后续工作是硬约束。

\subsection{第七步：推断的实现要点}

\begin{itemize}[leftmargin=1.6em]
  \item $\sqrt{\bar{s}_d}$ 在零点导数无穷，代码只在\emph{构造特征时}把谱截断到 $10^{-12}$。
        这保住了"严格零支撑"控制实验的语义，同时避免 NaN 梯度。
  \item 梯度裁剪到 $10$；$\log \sigma$ 与 $\log$ 方差都在对数空间优化以保证正性。
  \item routing 取 \texttt{global}（三个 mode 共享一组 $\boldsymbol{\pi}$）。
        逐 mode 或逐 rank 的 routing 在 $1\%$ 观测下过拟合，
        且对通用字典的伤害（$0.042$）大于对算子 bank（$0.023$），
        在那个设置下比较会虚高约 $1.6$ 倍。
\end{itemize}

\subsection{算法总览}

\begin{algorithm}[t]
\caption{构造逐维核 bank（\textbf{不使用任何观测数据}）}
\label{alg:bank}
\begin{algorithmic}[1]
\Require 算子族与系数 $\theta$（或范围 $\Theta$），网格 $\{n_d\}$，频率预算 $\{K_d\}$，atom 数 $Q$
\Ensure 每个 mode 一组非负谱 $\{s_{dq}\}$，本征基 $\{\Phi_d\}$
\For{每个空间 mode $d$}
  \State $\lambda_d[k] \gets (2 - 2\cos(\pi k / n_d)) / h_d^2$ \Comment{\cref{eq:eig}}
\EndFor
\State $\omega[j] \gets \pi j / (T \Delta t)$
\State $\Sop[j,a,b] \gets \bigl(\omega[j]^2 + (r + D_x \lambda_x[a] + D_y \lambda_y[b])^2\bigr)^{-1}$
       \Comment{\cref{eq:joint}}
\State $M \gets \mathbf{1}$；$M[0,\dots,0] \gets 0$ \Comment{掩码而非置零}
\State $\{s_{dq}\} \gets \textsc{NonnegCP}(\Sop, Q, M)$ \Comment{\cref{eq:cp}}
\State $s_{dq} \gets s_{dq} / \sum_k s_{dq}[k]$ \Comment{余弦基下的归一化}
\State $\Phi_d \gets [\,1,\ \sqrt{2}\cos(\pi k x)\,]$ \Comment{无流边界本征基}
\end{algorithmic}
\end{algorithm}

\begin{algorithm}[t]
\caption{拟合（\textbf{不使用任何留出数据}）}
\label{alg:fit}
\begin{algorithmic}[1]
\Require 观测 $(\Omega, \mathbf{y})$，bank $\{s_{dq}\}$，本征基 $\{\Phi_d\}$，秩 $\{R_d\}$，步数 $N$
\Ensure 变分参数 $\{\mathbf{a}^{(d)}\}$、核 $\mathcal{G}$、混合权重 $\pi$
\For{$\text{step} = 1 \dots N$}
  \State $\pi_d \gets \operatorname{softmax}(\text{logits})$；
         $\bar{s}_d \gets \sum_q \pi_{dq} s_{dq}$
  \State $\Psi_d \gets \Phi_d \odot \sqrt{\bar{s}_d}$ \Comment{Mercer 特征，自动半正定}
  \State $f^{(d)}_r(n) \gets \sum_k \Psi_d[i_d(n), k]\, \tilde{a}^{(d)}[r,k]$
         \Comment{重参数化采样}
  \State $\hat{u}(n) \gets \sum_{r_1 \cdots r_D} \mathcal{G}\prod_d f^{(d)}_{r_d}(n)$
         \Comment{只在 $\Omega$ 上求值}
  \State 最大化 ELBO，Adam 更新
\EndFor
\end{algorithmic}
\end{algorithm}

\paragraph{参数量。}主线设置下变分系数 $\sum_d R_d K_d = 216$（乘 $2$ 为均值与方差），
Tucker 核 $8 \cdot 5 \cdot 5 = 200$，\textbf{混合权重只有 $4$ 个}。
合计约 $636$ 个参数拟合 $2621$ 个观测。
"物理省下的是整个核，只留四个数要学"具体指的就是混合权重那四个。

%=====================================================================
\section{实验设定与数据}

\subsection{数据来源}

\textbf{全部主实验不依赖任何外部数据集}。场由本仓库的求解器现算：
DCT 谱方法加指数积分器处理扩散-反应（无流边界，余弦基对角化），
平流用算子分裂（谱步之后在实空间做一阶迎风，CFL 条件运行时检查）。
二维场约 $1$ 秒、三维场约 $1.4$ 秒一个 seed。

这样做是合理的，因为：真值来自\textbf{独立的}数值积分，
\emph{不是}从模型自己的先验里采样；而且给模型的名义系数是\textbf{故意错 $1.5$ 倍}的。

若要换真实数据，六个外部来源及其筛查结果见交接文档；其中
Kolmogorov flow（$\mathrm{Re}=40$）与 PDEBench 2D diffusion-reaction
已实测\textbf{不通过}低秩可行性筛查（各空间轴需 $63\%$ 的模态才够 $95\%$ 能量）。

\subsection{共享设定}

\begin{table}[t]
\centering
\caption{所有表共享的设定。同一 seed 内所有方法共享场、掩码与噪声，
因此每个比较都是配对的。}
\label{tab:settings}
\begin{tabular}{ll}
\toprule
真实场 & $64\times64$ 网格，$D=(0.02, 0.006)$，$r=0.04$，三处泄漏 \\
       & $\Delta t = 0.6$，burn-in $200$ 步，记录 $64$ 帧，背景噪声 $0.02$ \\
先验被告知 & $D = (0.03, 0.012)$，$r = 0.06$ \quad\textbf{（故意错 $50\%$）} \\
模型 & 频率数 $(12,12,12)$，Tucker 秩 $(8,5,5)$，atom 数 $4$，routing global \\
观测 & $1\%$（$2621$ 点），观测噪声 std $0.05$ \\
优化 & Adam，lr $0.02$，$1000$ 步，$3$ 样本 ELBO，梯度裁剪 $10$ \\
三维 & $32^4$（$105$ 万格点），频率数 $(8,8,8,8)$，秩 $(5,4,4,4)$，$400$ 步 \\
\bottomrule
\end{tabular}
\end{table}

\subsection{传感器布局}

所有布局的\textbf{观测预算完全相同}，只有排布不同：

\begin{table}[t]
\centering
\caption{传感器布局。$\mathrm{depth}$ 指到最近墙面的距离。}
\label{tab:layouts}
\begin{tabular}{lll}
\toprule
名称 & 区域 & 可达面积 \\
\midrule
散布 & 全域 & $100\%$ \\
四面墙 & $\mathrm{depth} < 2$ & $12\%$ \\
贴墙带 & $3 \le \mathrm{depth} < 8$ & $26\%$ \\
\textbf{单面墙} & $x < 5$ & $8\%$ \\
四个角块 & 四角各 $10\times10$ & $10\%$ \\
\midrule
\multicolumn{3}{l}{\emph{配对对照（用于隔离单一因素）}} \\
单面墙（$y$） & $y < 5$ & $8\%$ \\
左右两墙 / 上下两墙 & $x<5$ 或 $x \ge n_x-5$ / 同理 $y$ & $16\%$ \\
两面相邻墙 & $x<5$ 或 $y<5$ & $15\%$ \\
单个角块 & $x<20$ 且 $y<20$ & $10\%$ \\
\bottomrule
\end{tabular}
\end{table}

\paragraph{四个角块与单个角块的配对。}
这两个布局\textbf{观测格点数完全相同}（$400$ 格，$10\%$ 可达），
差别只在未观测区是被\emph{包围}还是仅\emph{相邻}。
把同样的面积摊成四个角，所有方法都变好，本方法从 $1.0346$ 降到 $0.8124$。
\textbf{这说明未观测区的形状比面积更决定难度。}
需要说明的是，这两个布局本方法都不占优（分别落后 $6.4\%$ 与 $2.6\%$），
它们的价值在于隔离出"形状 vs 面积"这一因素。

%=====================================================================
\section{结果}

\subsection{场景覆盖}

指标是留出集 $\NRMSE$（\cref{eq:nrmse}）。每格对比本方法与
\textbf{实践者能实际部署的最好那一臂}（在传感器读数上调参的 Matérn、
谱混合字典、神经张量补全，取三者最好）。

\begin{table}[t]
\centering
\caption{场景覆盖：三个算子族、两个维度、每种布局，$1\%$ 观测。}
\label{tab:coverage}
\small
\begin{tabular}{llccccc}
\toprule
 & & 散布 & 四面墙 & 贴墙带 & \textbf{单面墙} & 四个角块 \\
\midrule
\multicolumn{7}{l}{\emph{二维空间 $+$ 时间（$64^3$，$3$ seeds）}} \\
反应-扩散 & ours & $0.0547$ & $\mathbf{0.4766}$ & $\mathbf{0.3064}$ & $\mathbf{0.5387}$ & $0.8124$ \\
          & 最佳 baseline & $0.0530$ & $0.5004$ & $0.4206$ & $0.7719$ & $0.7917$ \\
          & 相对改进 & $-3.2\%$ & $+4.8\%$ & $\mathbf{+27.1\%}$ & $\mathbf{+30.2\%}$ & $-2.6\%$ \\
\addlinespace
扩散主导 & ours & $\mathbf{0.0388}$ & $\mathbf{0.2670}$ & $\mathbf{0.1909}$ & $\mathbf{0.3322}$ & --- \\
          & 最佳 baseline & $0.0395$ & $0.2895$ & $0.1947$ & $0.6720$ & --- \\
          & 相对改进 & $+1.7\%$ & $+7.8\%$ & $+1.9\%$ & $\mathbf{+50.6\%}$ & --- \\
\addlinespace
平流-扩散 & ours & $\mathbf{0.0265}$ & $0.4102$ & $\mathbf{0.2561}$ & $\mathbf{0.5144}$ & --- \\
          & 最佳 baseline & $0.0281$ & $0.3969$ & $0.3511$ & $0.7021$ & --- \\
          & 相对改进 & $+5.9\%$ & $-3.4\%$ & $\mathbf{+27.1\%}$ & $\mathbf{+26.7\%}$ & --- \\
\midrule
\multicolumn{7}{l}{\emph{三维空间 $+$ 时间（$32^4$，$2$ seeds；布局为面）}} \\
 & & 散布 & 对面两墙 & \textbf{单面墙} & 角落立方 & \\
反应-扩散 & ours & $0.1602$ & $\mathbf{0.5856}$ & $\mathbf{0.6856}$ & $\mathbf{0.9130}$ & \\
          & 最佳 baseline & $0.1586$ & $0.7356$ & $0.8883$ & $0.9733$ & \\
          & 相对改进 & $-1.1\%$ & $\mathbf{+20.4\%}$ & $\mathbf{+22.8\%}$ & $+6.2\%$ & \\
\bottomrule
\end{tabular}
\end{table}

\cref{tab:coverage}的读法：

\begin{itemize}[leftmargin=1.6em]
  \item \textbf{散布那一列，三个算子族、两个维度上全部打平}（$-3.2\%$ 到 $+5.9\%$）。
        这是应该的：随机采样下重构是内插，光滑先验已经够用。
  \item \textbf{单面墙那一列全部大幅改善}（$+22.8\%$ 到 $+50.6\%$）。
        这是外推，也正是本方法的适用区。
  \item 四面墙与贴墙带介于两者之间，改善随受限程度上升。
\end{itemize}

\subsection{与更强的对手比较}

\begin{table}[t]
\centering
\caption{与容量更大的模型、以及与同一份物理的另一种花法比较。
打星的臂\textbf{在留出区域上}选架构、学习率或残差权重，并有 $4$--$6$ 倍优化预算；
本方法一组配置、零调参数据。最后一列是同一网络关掉物理。}
\label{tab:baselines}
\small
\begin{tabular}{lcccccc}
\toprule
布局 & ours & PINN$^\star$ & CoSTCo$^\star$ & Fourier-MLP$^\star$ & LRTFR & 同网络无物理 \\
\midrule
散布 & $0.0547$ & $\mathbf{0.0522}$ & $0.1367$ & $0.1039$ & $0.0905$ & $0.1174$ \\
单面墙 & $\mathbf{0.5387}$ & $0.8008$ & $0.5924$ & $0.6060$ & $0.8966$ & $0.8229$ \\
\bottomrule
\end{tabular}
\end{table}

\paragraph{同一份物理，当先验比当残差惩罚值钱一个数量级。}
PINN 臂被告知的与本方法完全相同（方程形式 $+$ 错 $50\%$ 的系数），
都不知道泄漏源在哪。散布布局下残差买到 $0.065$ 并略胜本方法；
单面墙下同一个残差只买到 $0.022$，而且三个 seed 里有\textbf{两个}
在留出区域上的扫描直接选了残差权重 $\mathbf{0}$。

原因是结构性的：残差惩罚只在 collocation 点约束\emph{函数}，
而齐次方程有大量解，容量足够的网络能处处满足 $\Lop u \approx 0$
却在远离数据处全错；谱先验约束的是\emph{哪些函数先验上可能}——
这恰好是数据说完话之后还剩下的东西。

\subsection{需要多少方程知识}

\begin{table}[t]
\centering
\caption{知识档位（单面墙，$3$ seeds）。通用字典的 atom 数与池化 bank 完全相同，
且 collapsed 参数化保证变分系数量与 bank 大小无关，所以比较不被参数预算污染。}
\label{tab:ladder}
\begin{tabular}{lc}
\toprule
知道什么 & $\NRMSE$ \\
\midrule
K2：真实系数 & $0.5450$ \\
\textbf{K1：只知 $\theta^\star \in \times[1/3,\,3]$} & $\mathbf{0.5468}$ \\
K0：只知 $\theta^\star \in \times[1/10,\,10]$ & $0.5783$ \\
K$-$1：无物理（通用字典，atom 数相同） & $0.7408$ \\
\bottomrule
\end{tabular}
\end{table}

只知道系数落在 $\times[1/3, 3]$ 内，相对知道真值只差 $0.0018$，
相对无物理好 $0.194$。范围放宽到 $\times[1/10, 10]$ 后退化到 $0.5783$——
这正是本构造预先声明的可证伪预测：\textbf{bank 池化的范围足够宽时必然趋近通用字典}，
报告这个退化从哪里开始是主张的一部分，而不是尴尬。

%=====================================================================
\section{局限与尚未解释的现象}

\paragraph{本方法不赢调好的核，是打平。}
对着用留出区域选长度尺度的 Matérn，单面墙是 $0.5448$ 对 $0.5449$。
早期草稿曾声称 $+17.2\%$，那是错的：baseline 的长度尺度网格没有包含它自己的最优值。
\textbf{该结论已撤回。}

\paragraph{二维上随手固定一个合理常数就几乎追平。}
一个从不调参、直接固定 $\ell = 1.6$ 的实践者，在任何布局上相对本方法最差只差 $0.0135$。
本方法真正的价值集中在：跨场景时那个"合理常数"会变
（三维上需要的尺度跨 $20$--$30$ 倍），而方程自动给出它。

\paragraph{失败模式不体面。}
全部实验里唯一超过 $\NRMSE = 1.0$（比预测均值还差）的臂\textbf{始终是本方法}
（单个角块 $1.03$、$y$ 墙 $1.27$），任何 baseline 最差只到 $0.95$。
长尺度 Matérn 打不过时会收缩到均值并停下；本方法的尺度被方程钉死，
观测约束不住时会继续外推。\textbf{部署前需要一个检测"正在超出观测约束范围"
并转为收缩的回退机制，目前没有。}

\paragraph{两面墙的不对称尚未完全解释。}
同样 $8\%$ 可达，$x$ 墙 $0.5387$ 而 $y$ 墙 $1.2723$。
预先登记的"各向异性"解释被自己的各向同性对照推翻。
后续的剖面覆盖解释（贴 $x$ 墙的条带张满整个 $y$ 轴，
于是"只随 $y$ 变化"的分量被完全观测，而它占 $64\%$ 的方差）解释了约 $85\%$，
但交换扩散系数后优劣\emph{没有反转}，所以仍不完整。
需要指出：oracle Matérn 也表现出同样的不对称，
所以 $y$ 墙对\emph{所有方法}都更难——那是场的性质；
本方法额外退化得更厉害，才是本方法的问题。

\paragraph{标准物理信息 GP 在这个尺度上更快。}
AutoIP 式方法在二维规模上是 $1.7$--$8.0$ 秒，本方法是 $20$--$25$ 秒。
\textbf{不应写"它不 scale"}——在这个规模上该论断不成立，它的问题是精度。

%=====================================================================
\section{复现}

\begin{verbatim}
python experiments/run_leak_sensors.py --config base --tag leak_main3tier
python experiments/make_paper_tables.py
python experiments/make_leak_figures.py
\end{verbatim}

换一个场是改配置，不是改代码：

\begin{verbatim}
python experiments/run_leak_sensors.py --config advection_diffusion
python experiments/run_leak_sensors.py --config base --set evaluation.ratio=0.02
\end{verbatim}

未知键会报错而不是被忽略，展开后的完整配置写进每个结果 JSON，
所以任何数字都能追溯到产生它的设置。详见交接文档与快速上手文档。

\end{document}
